---
format: yohaku/0.2
target: text
macro: '[[macros/raw]]'
imported: {"bytes":4065,"path":"/root/projects/yohaku/release-prep/yohaku/agent/yohaku_text.py","sha256":"6220055947c73adfb2096545970d5aa5ec78224d90b043ffc1ab6e8ddf6ee5c7"}
---

## 正文
```text
#!/usr/bin/env python3
"""Human text-address interface for the same tools exposed to Yohaku's AI.

Examples (run from the Yohaku project directory):
  python3 agent/yohaku_text.py read files/demo.py.tex --view body
  python3 agent/yohaku_text.py replace files/demo.py.tex --view body \
    --start 1:8 --end 1:11 --revision SHA256 --expected old --text new

Use --text-file or --stdin for literal multiline replacements. No model/API
request is made. Paths are relative to the chosen world; all file operations
still use the Rust workspace and Loom.
"""
import argparse
import json
from pathlib import Path
import sys

import yohaku_agent as agent
from yohaku_coordinates import MAX_RANGE_BYTES


def point(value):
    try:
        parts = value.split(':')
        if len(parts) != 2:
            raise ValueError()
        line, column = map(int,parts)
        if min(line,column) < 1:
            raise ValueError()
        return line,column
    except ValueError:
        raise argparse.ArgumentTypeError('Use a positive LINE:COLUMN, for example 1:1') from None


def literal(path):
    with Path(path).open('rb') as stream:
        data = stream.read(MAX_RANGE_BYTES+1)
    if len(data) > MAX_RANGE_BYTES:
        raise ValueError('Replacement/expected text file exceeds 128 KiB')
    return data.decode('utf-8')


def main(argv=None):
    parser = argparse.ArgumentParser(description=__doc__,formatter_class=argparse.RawDescriptionHelpFormatter)
    parser.add_argument('--root',default=str(agent.ROOT/'worlds/default'))
    parser.add_argument('--loom',default=str(agent.ROOT/'bin/loom-linux-x86_64'))
    parser.add_argument('--config')
    commands = parser.add_subparsers(dest='operation',required=True)
    for operation in ('read','replace'):
        command = commands.add_parser(operation)
        command.add_argument('path',help='World-relative .tex or plain text path')
        command.add_argument('--view',choices=('body','text'),default='body')
        command.add_argument('--start',type=point,default=(1,1))
        command.add_argument('--end',type=point,required=operation=='replace')
        if operation == 'replace':
            command.add_argument('--revision',required=True)
            expected = command.add_mutually_exclusive_group(required=True)
            expected.add_argument('--expected')
            expected.add_argument('--expected-file')
            replacement = command.add_mutually_exclusive_group(required=True)
            replacement.add_argument('--text')
            replacement.add_argument('--text-file')
            replacement.add_argument('--stdin',action='store_true')
    args = parser.parse_args(argv)
    try:
        workspace = agent.Workspace(args.root,args.loom)
        config = agent.read_config(args.config)
        tools = agent.Tools(workspace.execute,permissions=config.get('permissions',agent.DEFAULT_PERMISSIONS),native=False)
        request = dict(path=args.path,view=args.view,start_line=args.start[0],start_column=args.start[1])
        if args.end:
            request.update(end_line=args.end[0],end_column=args.end[1])
        if args.operation == 'read':
            result = tools.call('read_range',request)
        else:
            expected = literal(args.expected_file) if args.expected_file else args.expected
            text = literal(args.text_file) if args.text_file else args.text
            if args.stdin:
                data = sys.stdin.buffer.read(MAX_RANGE_BYTES+1)
                if len(data) > MAX_RANGE_BYTES:
                    raise ValueError('Replacement text exceeds 128 KiB')
                text = data.decode('utf-8')
            request.update(revision=args.revision,expected=expected,text=text)
            result = tools.call('edit_range',request)
        print(json.dumps(result,ensure_ascii=False,indent=2))
        return 0 if result['ok'] else 1
    except (agent.ToolError,ValueError,OSError) as error:
        print(json.dumps({'ok':False,'error':str(error)},ensure_ascii=False))
        return 1


if __name__ == '__main__':
    sys.exit(main())
```

## 输出与关联
```yaml
output: "agent/yohaku_text.py"
```
